\section{Discussion}\label{sec:discussion}
The two robots expose different obstacles. On the CDPR, the isolated placement policy could handle an already-held object, yet the supervised fit under the full instruction tied its initializer at five completions in 128 attempts. Subsequent practice from empty-gripper starts gave the policy observations and actions at the actual transitions. On WidowX the problem appeared earlier. A policy that almost never picked up the object had few opportunities to practice carrying it after its own grasp. Improving its fit to stored actions left that bottleneck in place.

The CDPR comparison measures further training of a single policy lineage under a fixed milestone reward. The early terminal-only run used different resets, supervision, and success criteria. A causal comparison of the two rewards would require the same initializer, resets, termination rule, and interaction budget. Likewise, the LIBERO results compare fine-tuning configurations, with one reported endpoint per run. The observed 3.6-point gap at 160 updates does not establish a reproducible gain across training seeds.

All robot results are from simulation. The CDPR repeats measure rollout variation for two fixed checkpoints on development scenes used for model selection. Independent training runs, the reserved final-test scenes, and physical execution remain untested here. The contribution of scripted supervision also remains open. A direct test would train from the same skill policy with and without the assisted composition data, then compare empty-gripper success under one instruction.
